\SourceLeaf{303}%
points $A$, $B$, $C$, $D$ (I am using my own notation here, as I find \textsc{Newton}'s
rather cumbersome in this case) to find a point $\delta$ on that straight line.
He then uses $A$, $B$, $C$, $E$ to find another point $\varepsilon$ on the
line. The straight line itself is thus given in position, and for any other
point on it, $\mu$, his construction gives a corresponding
point $M$ on the conic. He lets the angle $BAC$ (with arms of unspecified length)
rotate about $A$ and notes the position in which $AB$ passes through $\mu$. Then $AC$ must
pass through $M$. Likewise, when the angle $ABC$ rotates about $B$, as
$BA$ passes through $\mu$, $BC$ must also pass through $M$. This construction is
exceedingly elegant, but one has command of the subject only when one can answer
all the other questions concerning this magical straight line.
In particular, one wants to know what relations this straight line has to the elements of the
conic, and whether those elements can themselves be derived easily from the position of the
straight line and the positions of $A$, $B$, $C$. I can now handle several things of this kind
quite elegantly, but I do not know whether I can carry the whole matter through myself,
since other duties force me to break off.

\mbox{.\,.\,.\,.\,.\,.\,.} Formerly, if I am not mistaken, you spent much time on conics.
You have certainly read much more on them than I have. For example, already in
Monatliche Correspondenz, 1810 November, p. 508, you cite an elegant proposition about the ellipse
from a book\HistoricalNote{[*)]}{[*)}{\textsc{Puissant}, Recueil de diverses propositions de géométrie. Paris 1801.
See also vol. IV, p. 396.]} that I have not yet consulted, although our library presumably
owns it. It is therefore natural to ask you first for historical information.
Do you know an elegant construction by which the \textls{center}
of a conic from 5 of its points --- \textls{without any unnecessary detours} --- may be
found? It certainly cannot be difficult, but this is precisely the derivation I have not yet
\mbox{worked out. .\,.\,.\,.\,.\,.\,.}
